% Owner: Phuc. This chapter describes methods present in Tuong/Bach's code.
\section{Observations, Labels and Analytical Scope}
The dataset describes web images with geometric and contextual features
\citep{kushmerick1999}. The processed labels are \texttt{ad.} and
\texttt{nonad.}. The models use \(Y=1\) for advertisements.

\section{Deduplication and Missing Values}
The duplicate key includes all original features and the label after
whitespace trimming. It excludes the exported row index and retains
identical feature vectors with conflicting labels. For each feature
\(j\), univariate summaries use the set \(O_j\) of records with observed
values. Class summaries also restrict this set to the relevant label.
The processed CSV preserves missing measurements and observed extremes.

\section{Numerical and Binary Summaries}
For observed values $x_1,\ldots,x_n$, the sample mean, corrected sample
variance and sample standard deviation are
\begin{equation}
\bar{x}=\frac{1}{n}\sum_{i=1}^n x_i,\qquad
s^2=\frac{1}{n-1}\sum_{i=1}^n(x_i-\bar{x})^2,\qquad s=\sqrt{s^2}.
\end{equation}
The standard deviation requires at least two observations. The median is the
middle ordered observation, or the average of the two middle observations
when $n$ is even. Mean and median summarize location; SD and the interquartile
range summarize spread \citep{nistHandbook}.

The EDA code uses R's type 7 sample quantiles. For probability $q$, define
$h=1+(n-1)q$, $j=\lfloor h\rfloor$ and $\gamma=h-j$. Between adjacent order
statistics the estimate is $(1-\gamma)x_{(j)}+\gamma x_{(j+1)}$; endpoints
return the minimum and maximum. The first and third quartiles use $q=0.25$
and $q=0.75$, and $\mathrm{IQR}=Q_3-Q_1$ \citep{rQuantile}.

For a binary feature and class $c$, the observed prevalence is
\begin{equation}
\hat{p}_{jc}=\frac{\#\{i\in O_j:y_i=c,\ x_{ij}=1\}}
{\#\{i\in O_j:y_i=c\}}.
\end{equation}
Its denominator excludes missing values. Class proportions instead use the
total number of labeled records. Differences between binary-feature rates
are expressed in percentage points, $100(\hat p_{j,\mathrm{ad}}-
\hat p_{j,\mathrm{nonad}})$. Ranking these differences across many indicators
is exploratory screening, not a significance test or validated feature selection.

\section{Graphical Exploration and Extreme Values}
Histograms display the shape of observed distributions. Class-specific
boxplots compare medians, spreads and potential extremes; scatterplots reveal
overlap and relationships between numerical measurements. These plots complement the numerical summaries \citep{nistHandbook}.

The numerical outlier audit uses pooled type 7 quartiles and flags values
strictly below $Q_1-1.5\,\mathrm{IQR}$ or above $Q_3+1.5\,\mathrm{IQR}$.
Flagging identifies observations for inspection; it does not justify deletion.
The boxplots use R's default Tukey hinges, which need not coincide exactly
with type 7 quartiles \citep{rBoxplot}. Therefore their displayed fliers may
differ from the pooled audit counts, especially when plots are split by class.
All observed extremes are retained in the current workflow.

A normal Q--Q plot compares ordered observations with theoretical normal
quantiles. The implemented reference line passes through the first and third
quartiles. Departures from a straight line indicate limitations of a normal
approximation to that variable \citep{rQQ}. This graphical assessment is not
a formal normality test, nor does it determine the assumptions of a future
classification model.

\section{Pearson Correlation and Complete Cases}
For paired observations with nonzero variation in each variable, Pearson's
sample correlation is
\begin{equation}
r_{XY}=\frac{\sum_{i=1}^{n}(x_i-\bar{x})(y_i-\bar{y})}
{\sqrt{\sum_{i=1}^{n}(x_i-\bar{x})^2\sum_{i=1}^{n}(y_i-\bar{y})^2}}.
\end{equation}
It measures the direction and strength of linear association. The three correlations use the same records with complete geometry.
Aspect ratio is derived from dimensions, so some association is expected.
Complete records can have a different class balance from the full data.


\section{Comparing Means and Testing Associations}
Welch's two-sample $t$-test compares independent means while allowing the
group variances to differ. Its standard error is
$\sqrt{s_1^2/n_1+s_2^2/n_2}$ and its degrees of freedom use the
Welch--Satterthwaite approximation \citep{rWelch}. We apply the test
to observed image widths in the training set.

Pearson's chi-square statistic for a contingency table is
\(\chi^2=\sum_{ij}(O_{ij}-E_{ij})^2/E_{ij}\), where
\(E_{ij}=n_{i\cdot}n_{\cdot j}/n\) under independence. For a two-by-two
table, the test uses one degree of freedom and Cramer's
\(V=\sqrt{\chi^2/n}\) measures association size. Chapter~5 uses the
uncorrected statistic and checks that expected counts are at least five.
Term-level tests use the Benjamini--Hochberg procedure across eligible
tests to address false discoveries under its applicable assumptions
\citep{benjamini1995}. The results remain exploratory because terms can
be correlated or share identical occurrence patterns.

\section{Logistic Regression and Nonlinear Geometry}
With \(Y=1\) denoting an advertisement, a binomial logistic model relates
the probability \(p=\Pr(Y=1\mid X)\) to predictors through
\(\log[p/(1-p)]=\beta_0+\sum_j\beta_jX_j\)
\citep{rGLM}. In the main effects model, a unit increase in a predictor
multiplies conditional odds by \(\exp(\beta_j)\). A unit increase in
\(\log_2\) width or height is a doubling of the original dimension.
An approximate 95\% Wald interval is obtained by exponentiating
\(\hat\beta_j\pm1.96\,\mathrm{SE}(\hat\beta_j)\).
The likelihood ratio test for the added terms compares the drop in
deviance of the nested models against a chi-square approximation.
The tests assume independent observations and do not adjust for screening.

The extended geometry model uses natural cubic spline bases with three
degrees of freedom for each log dimension, a wide-banner indicator, and
the product of the log dimensions. Splines allow curved effects on
log odds and become linear beyond the boundary knots \citep{rSplines}.
The interaction permits the change in log odds associated with one
dimension to depend on the other. The modelling code estimates missing
dimension medians and the mode of Local from each fitting sample.
The processed CSV used for descriptive analysis keeps the original
missing values.

\section{Grouped Validation and Performance Measures}
Groups defined by identical feature vectors are kept together in the
holdout and the five training folds. Each out-of-fold (OOF) score comes
from a model fitted without the record's fold. Imputation and term
selection are repeated inside each fold. Thresholds maximize OOF recall
among cutoffs reaching precision 0.70 on training predictions and are
then held fixed on the test set.

With advertisements as positives, precision is \(TP/(TP+FP)\) and
recall is \(TP/(TP+FN)\). The precision--recall curve shows the tradeoff across thresholds and is
particularly relevant when positives are uncommon \citep{davis2006}.
Average precision sums \(P_k(R_k-R_{k-1})\) over distinct score
cutoffs, with tied scores admitted together. ROC-AUC measures ranking
across positive and negative records, with half credit for tied scores.
Neither ranking measure certifies performance at a specific cutoff.


\section{Uncertainty in Predictive Comparisons}
Wilson intervals invert the binomial score test. In the test set the
precision denominator is the number flagged, and the recall denominator
is the number of ads \citep{rWilson}. The exact McNemar test compares
paired decisions using a two-sided binomial test with probability $1/2$
on the two discordant counts \citep{rBinomial}. We compare ads and
non-ads separately to distinguish detections from false alarms.

The paired bootstrap resamples test records within each class, keeping
their two scores together \citep{efron1979}. The 2.5th and 97.5th percentiles
of 2,000 AP differences give an interval conditional on this split and fit.
